\scanpage{233}
% composite defining the map from the preceding leaf
est l'hom. composé [linéaire pour la structure \uline{gauche} sur $\unsure{P|p}^{i}(M)$]
\diagram{\begin{tikzcd}[ampersand replacement=\&]
 E\arrow[r,"{\delta_i}"]\&p^i(E)\arrow[r]\&p^i(M).
\end{tikzcd}}

% representing module and inverse system
Donc, pour $P$ variable, c'est représenté par
\[
 \Phi_i=p^i(M)/\underline{\mathcal O}_S\cdot\rho_i(E),
\]
les $\Phi_i$ \unsure{formant|forment} un syst. projectif strict, déterminé par la connaissance de $M$ et de
\[
 \rho_\infty:E\longrightarrow P^\infty(M)
\]
(linéaire pour la structure \uline{gauche} de $P^\infty(M)$), défini par les $\rho_i$. Attention, \unsure{ce|ici} $\rho_\infty$ satisfait à des \deleted{\unsure{deux}} conditions spéciales : son composé $E\longrightarrow M$ avec l'augmentation $P^\infty(M)\longrightarrow M$ est un épimorphisme, $+$ conditions d'intégrabilité---
% archive sheet number: 232
